\documentclass[11pt]{article}

\usepackage{amsmath, amssymb, amsthm, geometry}
\usepackage{tikz-cd}
\usepackage{booktabs}
\usepackage{hyperref}
\usepackage{array}
\usepackage{calc}
\usepackage{color}
\geometry{margin=1in}

\newtheorem{theorem}{Theorem}[section]
\newtheorem{proposition}[theorem]{Proposition}
\newtheorem{lemma}[theorem]{Lemma}
\newtheorem{definition}[theorem]{Definition}
\newtheorem{remark}[theorem]{Remark}
\newtheorem{conjecture}[theorem]{Conjecture}
\newtheorem{corollary}[theorem]{Corollary}

\providecommand{\tightlist}{%
  \setlength{\itemsep}{0pt}\setlength{\parskip}{0pt}}

\begin{document}

\title{L-S Contraction and the Boundary of Applicability:\\
$H^3$, the 16-Cell, and Non-Associative Algebra\\
{\large (A Dynamical Note toward $H^3$ (Borromean Contextuality / Bundle Gerbes))}}

\author{Zhou-Li Chen \\ co-nlang Research}
\date{May 2026}

\maketitle

\begin{abstract}
This note ascends the obstruction ladder to $H^3$, examining both the
16-cell nerve construction of Borromean contextuality and the fundamental
limit posed by non-associative algebras. In the associative setting
(Pauli group on 5 qubits), the 3-simplex overlaps of the 16-cell produce
an associator anomaly classified by the $d_3$ differential in the LHS
spectral sequence---the dynamical analogue of a bundle gerbe
Dixmier--Douady class. Beyond this, when the underlying algebra is
non-associative (octonions $\mathbb{O}$), the L-S contraction framework
itself ceases to be well-defined: the chain rule fails, the Jacobian
is not unique, and the concept of ``continuous deterministic dynamics''
loses meaning. This provides a purely dynamical proof that
$\mathbb{O}$-valued quantum mechanics cannot support stable classical
observation.
\end{abstract}

%------------------------------------------------------------------
\section{From 4-Cycles to 3-Simplices}
\label{sec:intro}
%------------------------------------------------------------------

The obstruction ladder ascends through increasing dimension of the
MASA nerve:

\[
\begin{array}{cccc}
\text{Level} & \text{Nerve structure} & \text{L-S obstruction} & \text{Cohomology} \\
\hline
H^1 & \text{1-cycle (loop)} & \det\Phi_\gamma < 1 \text{ vs } \det\Phi_\gamma = 1 & \check{H}^1(M, U(1)) \\
H^2 & \text{4-cycle ($K_{3,3}$)} & \text{Central holonomy } \omega = -1 & H^2(G/N, \{\pm 1\}) \\
H^3 & \text{3-simplex (tetrahedron)} & \text{Associator anomaly } \Delta \neq 0 & H^3(M, \mathbb{Z}) \text{ (DD class)} \\
\end{array}
\]

At $H^1$, a single loop of MASAs fails to close (Liouville vs.\ contraction).
At $H^2$, a 4-cycle of pairwise-compatible transitions accumulates a central
sign. At $H^3$, four mutually overlapping MASAs on a $2$-sphere nerve produce
an associator anomaly: the composition of three L-S transition matrices
depends on the order in which they are grouped.

\begin{remark}
In the associative setting (Pauli group), matrix multiplication is strictly
associative. The $H^3$ obstruction does NOT arise from non-associativity of
the operators themselves, but from the \textbf{higher categorical} structure
of context-switching. Two different ways of patching three transitions
around a 3-simplex differ by a $U(1)$ phase---a 3-cocycle in the
\v{C}ech nerve. This is the dynamical analogue of the Dixmier--Douady
class of a bundle gerbe~\cite{Murray96,Brylinski93}.
\end{remark}

%------------------------------------------------------------------
\section{The 16-Cell Nerve and Its 3-Simplices}
\label{sec:16cell}
%------------------------------------------------------------------

Paper IV~\cite{PaperIV} established that the minimal configuration for
Borromean contextuality is the 16-cell (4-dimensional cross-polytope):
8 MASA vertices forming 4 antipodal pairs, with the homotopy type of
$S^3$. The 16-cell contains $2^4 = 16$ tetrahedra (3-simplices),
each corresponding to a quadruple MASA overlap.

For $N=5$ qubits (the minimal capacity to avoid symplectic collapse),
the 16 core operators $\{P_{b_3b_2b_1b_0}\}$ satisfy:

\[
[P_x, P_y] = 0 \quad \text{unless } y = \sim x \text{ (bitwise complement)}.
\]

Each MASA $M_{d,b}$ ($d \in \{0,1,2,3\}$, $b \in \{L,R\}$) contains the
8 core operators whose $d$-th index bit equals $b$.

\subsection{Transition matrices on tetrahedra}
\label{sec:transitions}

For a 3-simplex (tetrahedron) formed by four MASAs
$M_1, M_2, M_3, M_4$ that are pairwise compatible (each pair shares a
core operator), define the L-S transition matrix
$\mathbf{T}_{ij}: \mathbb{C}^{32} \to \mathbb{C}^{32}$ as the unitary
change-of-basis between the eigenbases of $M_i$ and $M_j$. The
$32$-dimensional Hilbert space reflects $N=5$ qubits.

\begin{lemma}[3-simplex transition composition]
\label{lem:3simplex-comp}
Let $\mathbf{T}_{12}, \mathbf{T}_{23}, \mathbf{T}_{34}$ be L-S transition
matrices around three edges of a tetrahedron. The two ways of composing
them differ by a $U(1)$ phase:
\[
(\mathbf{T}_{12}\mathbf{T}_{23})\mathbf{T}_{34}
= e^{i\theta} \; \mathbf{T}_{12}(\mathbf{T}_{23}\mathbf{T}_{34}).
\]
In the 5-qubit Pauli group, $e^{i\theta} = \pm 1$.
\end{lemma}

\begin{proof}[Sketch]
Each $\mathbf{T}_{ij}$ is an element of the 5-qubit Clifford group (up to
phase), which is a subgroup of the Pauli group $\mathcal{P}_5^H$ modulo
central phases. The composition of three Cliffords around a tetrahedron
edge returns an element of $\mathcal{P}_5^H$; the associator
\[
\Delta = (\mathbf{T}_{12}\mathbf{T}_{23})\mathbf{T}_{34}
- \mathbf{T}_{12}(\mathbf{T}_{23}\mathbf{T}_{34})
\]
measures the difference between two ways of grouping the three transitions.
Because the two groupings differ only by the placement of parentheses,
and each $\mathbf{T}_{ij}$ is an element of the 5-qubit Clifford group
(whose products land in $\mathcal{P}_5^H$), the difference $\Delta$ is a
central element $\Delta = (\alpha - 1)\,\mathbf{I}$ for some phase
$\alpha \in U(1)$. This phase is a 3-cochain on the nerve of $\mathcal{C}$;
the LHS spectral sequence's $d_3$ differential maps this 3-cochain to the
generator of $H^3(\bar{\mathcal{P}}_5, U(1)) \cong \mathbb{Z}/2$
when $\alpha \neq 1$~\cite{PaperIV}.
\end{proof}

%------------------------------------------------------------------
\section{The $d_3$ Transgression}
\label{sec:d3}
%------------------------------------------------------------------

In the LHS spectral sequence for the 5-qubit central extension, the
$d_3$ differential maps:

\[
d_3: E_3^{0,2} \longrightarrow E_3^{3,0}.
\]

The domain $E_3^{0,2}$ consists of classes that survive the $d_2$
transgression---those 2-cohomology obstructions that are not eliminated
by the central extension. Their image under $d_3$ lies in
$H^3(\bar{\mathcal{P}}_5, U(1))$, which for the 5-qubit Pauli group is
non-trivial~\cite{PaperIV}.

\begin{remark}[Geometric interpretation]
The $d_3$ differential measures the associativity of the L-S transition
composition on 3-simplex overlaps. Just as $d_2$ captured the central
sign flip around a 4-cycle (non-commutativity), $d_3$ captures the
``phase of a phase'' around a 3-simplex (non-associativity of
context-patching).
\end{remark}

%------------------------------------------------------------------
\section{The $H^3$ Correspondence: Global Conjecture}
\label{sec:conjecture}
%------------------------------------------------------------------

The 16-cell nerve $|N\mathcal{C}_{16}|$ is homotopy equivalent to $S^3$,
with $\check{H}^3(|N\mathcal{C}_{16}|, \mathbb{Z}) \cong \mathbb{Z} \neq 0$.
Unlike the $H^2$ case---where a single 4-cycle computation sufficed---the
$H^3$ obstruction is intrinsically \textbf{global}: it is the \v{C}ech
3-cocycle on the entire $S^3$ nerve, not on any single 3-simplex.

\begin{conjecture}[Global $S^3$ Dynamical Obstruction]
\label{conj:ls-h3}

\noindent\textbf{(1) Geometric premise.}
The 16-cell nerve $|N\mathcal{C}_{16}|$ of the 5-qubit Borromean system
satisfies $|N\mathcal{C}_{16}| \simeq S^3$, hence
$\check{H}^3(|N\mathcal{C}_{16}|, \mathbb{Z}) \cong \mathbb{Z} \neq 0$.

\noindent\textbf{(2) Global coboundary of face 2-cochains.}
The L-S transition matrices $\{\mathbf{T}_{ij}\}$ on the edges of the
16-cell define, for each of the 24 triangular faces, a \v{C}ech
2-cochain $c_{ijk} \in \mathbb{Z}/2$ (the face holonomy, analogous to
the $H^2$ obstruction). The total \v{C}ech 3-cocycle
$[c_3] \in \check{H}^3(|N\mathcal{C}_{16}|, \underline{U(1)})$ is the
coboundary of $\{c_{ijk}\}$ over the 16 tetrahedra. This 3-cocycle
measures the obstruction to a global L-S contraction metric on the
$S^3$ nerve.

\noindent\textbf{(3) The dynamical pairing hypothesis.}
By analogy with the $H^2$ case (where the central $-\mathbf{I}$ arises
from a 4-cycle holonomy), we conjecture that the pairing of $[c_3]$
with the fundamental class $[S^3]$ is non-trivial:
\[
\langle [c_3], [S^3] \rangle = -1
\quad\in\quad H_3(S^3, \mathbb{Z}) \otimes \check{H}^3(S^3, \underline{U(1)}).
\]
This $-1$ represents the \textbf{global sign flip} that a continuous
classical contraction metric incurs when forced to wrap the entire
observation universe ($S^3$).

Consequently, Borromean contextuality corresponds to a non-trivial
$d_3$ transgression in the LHS spectral sequence, and the
Dixmier--Douady class $\mathrm{DD}(\mathcal{G}) \in H^3(M, \mathbb{Z})$
is the topological charge of the L-S contraction obstruction on the
global $S^3$ nerve.
\end{conjecture}

\begin{remark}[Computational horizon]
\label{rem:computational}
The 16-cell contains 8 vertices, 32 edges, 24 triangular faces, and
16 tetrahedra. Verifying that the coboundary of the 24 face 2-cochains
yields a non-trivial 3-cocycle with $\langle [c_3], [S^3] \rangle = -1$
requires a SAT solver or automated theorem prover (see Paper IV's
supplementary material~\cite{PaperIV}). This is identified as an open
computational problem for future work.
\end{remark}

%------------------------------------------------------------------
\section{Problem B: The Octonionic Boundary of L-S Theory}
\label{sec:octonionic-boundary}
%------------------------------------------------------------------

The obstruction ladder ($H^1 \to H^2 \to H^3$) is built entirely within
associative algebras ($\mathbb{C}$-valued or $\mathbb{H}$-valued
representation theory). However, the algebraic counterpart of $H^3$
is the exceptional Jordan algebra $\mathfrak{h}_3(\mathbb{O})$---the
Albert algebra, which is inherently non-associative.

Here L-S theory encounters a \textbf{foundational limitation}, not
merely a cohomological obstruction.

\subsection{The chain rule is the issue, not associativity per se}
\label{sec:chain-rule}

The L-S contraction analysis requires three structures:
\begin{enumerate}
  \item A smooth flow $\dot{\mathbf{z}} = \mathbf{f}(\mathbf{z})$ on a manifold.
  \item The variational equation $\delta\dot{\mathbf{z}} = \mathbf{J}\,\delta\mathbf{z}$,
    where $\mathbf{J}_{ij} = \partial f_i / \partial z_j$.
  \item A symmetric positive-definite metric $\mathbf{M}$ that makes
    $\mathbf{M}\mathbf{J} + \mathbf{J}^T\mathbf{M} + \dot{\mathbf{M}}$ negative-definite.
\end{enumerate}

Step 2 is the critical point. The Jacobian $\mathbf{J}$ is well-defined
because the derivative $\partial/\partial z_j$ is a linear operator
satisfying the chain rule:
\[
\frac{\partial}{\partial z_j}(g \circ h) = \sum_k 
\frac{\partial g}{\partial h_k}\,\frac{\partial h_k}{\partial z_j}.
\]
If the state space $\mathcal{M}$ is an algebra over $\mathbb{O}$, and if
the dynamics $\mathbf{f}$ involves octonionic multiplication (e.g.,
$\dot{x} = ax$ for $a, x \in \mathbb{O}$), then repeated differentiation
produces expressions where the chain rule is ambiguous:
\[
\frac{\partial}{\partial t}(x \cdot y) = \dot{x} \cdot y + x \cdot \dot{y}
\quad\text{(holds for $\mathbb{O}$)},
\]
but
\[
\frac{\partial}{\partial t}((x \cdot y) \cdot z) \neq
\frac{\partial}{\partial t}(x \cdot (y \cdot z))
\quad\text{(because $(xy)z \neq x(yz)$)}.
\]
Consequently, the Jacobian matrix $\mathbf{J}$ is \textbf{not uniquely
defined}: its components depend on the order in which parentheses are
placed in the expression $\partial f_i / \partial z_j$.

\begin{theorem}[L-S theory fails for $\mathbb{O}$-dynamics]
\label{thm:o-failure}
Let $\dot{\mathbf{z}} = \mathbf{f}(\mathbf{z})$ be a dynamical system on
$\mathbb{O}^n$ whose flow involves non-associative multiplication.
Then the variational equation $\delta\dot{\mathbf{z}} = \mathbf{J}\,\delta\mathbf{z}$
is not uniquely defined: the Jacobian $\mathbf{J}$ depends on a choice
of parenthesization and is therefore not a well-defined linear operator.
Consequently, L-S contraction analysis, and with it the concept of
``stable classical observability,'' cannot be formulated.
\end{theorem}

\begin{proof}
The dynamics $\dot{z}_i = f_i(z_1, \dots, z_n)$ contains at least one
term involving a product of two or more state variables. Since
$(ab)c \neq a(bc)$ in $\mathbb{O}$, the partial derivative
$\partial f_i / \partial z_j$ depends on whether the term $f_i$ is
computed as $(z_j z_k)z_l$ or $z_j(z_k z_l)$. These two evaluations
give different results for $\delta\dot{z}_i$ when $\delta z_j \neq 0$,
so $\mathbf{J}$ is not a single-valued linear map.
\end{proof}

\subsection{Physical interpretation}
\label{sec:physical}

The failure of the Jacobian to be well-defined has a direct physical
meaning:
\begin{itemize}
  \item \textbf{$\dot{\mathbf{z}} = \mathbf{f}(\mathbf{z})$} describes the
    system's deterministic evolution (kinematics). This survives in
    $\mathbb{O}$: one can solve $\ddot{x} = -x$ on $\mathbb{O}$.
  \item \textbf{$\delta\dot{\mathbf{z}} = \mathbf{J}\,\delta\mathbf{z}$}
    describes the propagation of infinitesimal perturbations---the
    basis of prediction, stability analysis, and measurement theory.
    This \emph{fails} in $\mathbb{O}$ because the chain rule is
    ambiguous.
\end{itemize}
The conclusion: in an $\mathbb{O}$-valued quantum mechanics,
\textbf{macroscopic classical observability (L-S contraction) cannot
emerge}, because the fundamental tool that guarantees the uniqueness
of perturbation propagation---the chain rule---is absent. This provides
a purely \textbf{dynamical proof} that the underlying division ring
of physics must be associative.

\begin{remark}[Implication for $\mathbb{O}$-valued physics]
\label{rem:ls-H3-implication}
Theorem~\ref{thm:o-failure} shows that L-S contraction analysis---the only
known dynamical framework guaranteeing unique perturbation propagation---cannot
be formulated in $\mathbb{O}$. This does not prove that \emph{no} framework for
stable classical observation exists in $\mathbb{O}$, but it removes the most
powerful tool available. Together with the $d_2$ argument of
Paper VI~\cite{PaperVI} (which selects $\mathbb{C}$ from the associative
division algebras), the status of $\mathbb{O}$ as a physical amplitude algebra
is doubly excluded: by algebraic necessity (Paper VI) and by dynamical
inapplicability (Theorem~\ref{thm:o-failure}).
\end{remark}

%------------------------------------------------------------------
\section{Conclusion: The Dynamical Origin of the Obstruction Ladder}
\label{sec:conclusion}
%------------------------------------------------------------------

This note completes the trilogy of L-S dynamical appendices to the
cohomological obstruction ladder. The central result is a purely
dynamical reinterpretation of quantum contextuality:

\begin{center}
\begin{tabular}{llll}
\toprule
Level & L-S obstruction & Cohomology & Status \\
\midrule
$H^1$ & Liouville $\det\Phi = 1$ vs.\ contraction $\det\Phi < 1$ &
$\check{H}^1(M, \underline{U(1)})$ & \textbf{Theorem} \\
$H^2$ & 4-cycle central holonomy $\omega = -1$ &
$H^2(G/N, \{\pm 1\})$ & \textbf{Theorem} \\
$H^3$ & Global 3-cocycle $\langle [c_3], [S^3] \rangle \stackrel{?}{=} -1$ &
$\check{H}^3(M, \underline{U(1)})$ & \textbf{Conjecture} \\
$\mathbb{O}$ & Chain rule fails; L-S Jacobian not well-defined &
--- & \textbf{Theorem} \\
\bottomrule
\end{tabular}
\end{center}

The four levels are not merely a taxonomic ladder. They trace the
progressive failure of increasingly foundational structures in classical
deterministic dynamics when forced onto the non-classical geometry of a
quantum state space:

\begin{enumerate}
  \item \textbf{Topological} ($H^1$): A loop of classical measurement
    contexts fails to close because Liouville forbids global contraction.
    $\to$ Geometric phase.
    \hfill\textit{[Proven in companion appendix.]}

  \item \textbf{Group-theoretic} ($H^2$): Non-commuting transition matrices
    around a 4-cycle of the $K_{3,3}$ nerve accumulate a central sign
    $-\mathbf{I}$. $\to$ Kochen--Specker contextuality.
    \hfill\textit{[Proven in companion note.]}

  \item \textbf{Categorical} ($H^3$): The 16-cell nerve $S^3$ carries a
    global Čech 3-cocycle. Its pairing with the fundamental class is
    conjectured to be $-1$, representing the obstruction to a globally
    contracting classical metric on the full observation space.
    $\to$ Borromean contextuality.
    \hfill\textit{[Conjecture; SAT verification open.]}

  \item \textbf{Dynamical} ($\mathbb{O}$): In a non-associative algebra,
    the chain rule fails and the Jacobian is not uniquely defined.
    L-S contraction theory is inapplicable at the axiomatic level.
    $\to$ No stable classical observation possible.
    \hfill\textit{[Theorem~\ref{thm:o-failure}]}
\end{enumerate}

The common thread across all four levels is a single insight:
\textbf{quantum contextuality is the dynamical memory of classical
systems that refuse to contract.} The L-S framework translates the
abstract cohomology of MASA posets into concrete statements about
Liouville volume preservation, symplectic geometry, unitary transition
holonomy, and the limits of the chain rule. The obstruction ladder is
not built from pure algebra---it is carved from the impossibility of
stable, forgetting dynamics in a quantum universe.

\begin{thebibliography}{9}

\bibitem[Murray96]{Murray96}
M.~K.~Murray,
``Bundle gerbes,''
\textit{J. London Math. Soc.} \textbf{54} (1996), 403--416.

\bibitem[Brylinski93]{Brylinski93}
J.-L.~Brylinski,
\textit{Loop Spaces, Characteristic Classes and Geometric Quantization}.
Progress in Mathematics, Vol.~107, Birkh\"auser, 1993.

\bibitem[PaperIV]{PaperIV}
Chen, Z.-L.
The Cohomological Obstruction Ladder.
\textit{Zenodo, DOI: 10.5281/zenodo.20073184}, 2026.

\bibitem[PaperVI]{PaperVI}
Chen, Z.-L.
The Ultimate Axiom: Deriving Quantum Mechanics from the Logic of
Observation.
\textit{Zenodo, DOI: 10.5281/zenodo.20073190}, 2026.

\end{thebibliography}

\end{document}
